\section{局部一阶谱几何}

本节说明为何 polar/SpecGrad 更新自然对应 operator-norm 几何，而非 Frobenius 几何。

\subsection{SpecGrad 家族}

设
\[
G=U\Sigma V^\top
\]
为某个矩阵梯度。定义谱梯度更新族：
\[
D_f(G)=Uf(\Sigma)V^\top.
\]
特殊情形包括：
\[
f(\sigma)=\sigma \quad \Rightarrow \quad D_f(G)=G,
\]
\[
D_F(G)=G/\|G\|_F,
\]
以及 polar 方向
\[
f(\sigma)=1_{\sigma>0},\qquad D_{\mathrm{polar}}(G)=UV^\top.
\]

\subsection{Frobenius 与 operator-norm 约束下的一阶最优性}

在当前参数点，对单个矩阵块考虑局部展开：
\[
\Loss(W-D)-\Loss(W)
= -\langle G,D\rangle+O(\|D\|^2).
\]
因此，一阶目标是最大化
\[
\langle G,D\rangle.
\]

\begin{proposition}[谱分配边界]
设 \(G=U\operatorname{diag}(\sigma)V^\top\)。
\begin{enumerate}[label=(\roman*)]
  \item 在 Frobenius 球 \(\|D\|_F\le r\) 上，唯一最优方向为
  \[
  D_F^\star=rG/\|G\|_F,
  \]
  最优一阶下降为
  \[
  r\|G\|_F.
  \]
  \item 在 operator-norm 球 \(\|D\|_{\op}\le \eta\) 上，一个最优方向为
  \[
  D_{\op}^\star=\eta UV^\top,
  \]
  最优一阶下降为
  \[
  \eta\|G\|_*.
  \]
\end{enumerate}
\end{proposition}

\begin{proof}
Frobenius 情况由 Cauchy--Schwarz 给出：
\[
\langle G,D\rangle\le \|G\|_F\|D\|_F\le r\|G\|_F,
\]
等号在 \(D=rG/\|G\|_F\) 处取得。

Operator-norm 情况由 von Neumann trace inequality 给出：
\[
\langle G,D\rangle
\le \sum_i\sigma_i(G)\sigma_i(D)
\le \eta\sum_i\sigma_i(G)
=\eta\|G\|_*.
\]
取 \(D=\eta UV^\top\) 达到等号。
\end{proof}

该命题解释了一个关键事实：
\[
\boxed{\text{polar/flat spectrum 不是普适最优，而是 operator-norm 预算下的一阶最优。}}
\]
如果公平比较的更新预算是 Frobenius 范数，那么梯度形状 \(G/\|G\|_F\) 才是一阶最优；如果公平比较的是 operator norm，那么 polar 更新才是一阶最优。
